\documentclass[12pt]{article}
\usepackage[utf8]{inputenc}
\usepackage[T1]{fontenc}

\title{Taller: Organización y Búsqueda de Registros}
\author{}
\date{}

\begin{document}

\maketitle

Cuenta la historia que en una antigua biblioteca escondida bajo las arenas del desierto,
una dinastía de guardianes custodiaba y mantenía un archivo milenario con los nombres
y destinos de todas las personas que han existido y existirán. Estos nombres debían ser
organizados de tal forma que, sin importar cuántos nombres se fueran agregando, siempre
pudiera encontrarse el archivo de cualquier persona conociendo únicamente su nombre.

Sin embargo, con el paso de las generaciones, el conocimiento sobre cómo se debían
agregar los nombres de manera que fuera fácil consultarlos se fue perdiendo, generando un
registro caótico de 100.000 nombres que deben ser reubicados dentro del archivo.

La generación actual de guardianes realizó un minucioso inventario de estos nombres, pero,
debido a que no pudieron ponerse de acuerdo entre sí, actualmente cuentan con tres listas
diferentes:

\begin{itemize}
    \item Una lista con los nombres en el orden en que fueron registrados por los escribas (orden aleatorio).
    \item Una lista con los nombres ordenados alfabéticamente.
    \item Una lista con los nombres en orden alfabético inverso.
\end{itemize}

Los guardianes llegaron a la conclusión de que su problema se dividía en dos problemas
más pequeños:

\begin{itemize}
    \item ¿Cómo podrían insertar los nombres dentro del archivo de la manera más eficiente posible, sin importar el orden en el que decidieron organizar los nombres?
    \item Una vez insertados los nombres dentro del archivo, ¿cuál de las estructuras que usaron como base para la inserción permite una búsqueda más rápida de los nombres dentro del archivo?
\end{itemize}

Para responder a esta pregunta se deben satisfacer los siguientes requerimientos:

\begin{itemize}
    \item Probar cuál de las estructuras de datos jerárquicas vistas en clase hasta el momento (BST, AVL, RBT) se muestra más efectiva para solucionar el problema.
    \item Evaluar los resultados de los tiempos de ejecución para inserción y posteriormente para la búsqueda de hasta 25 registros seleccionados aleatoriamente.
    \item Analizar cómo influye el orden de inserción en el rendimiento de cada estructura.
\end{itemize}

\textbf{Pista:} Esta tarea puede realizarse con un arreglo de nombres generados aleatoriamente como insumo. Un nombre es una cadena con un tamaño de no menos de 3 y no más de 10 caracteres. Si estos nombres son almacenados en un arreglo de manera previa, la tarea de buscar un registro aleatorio es más fácil de realizar.

De manera adicional, se sugiere realizar varias ejecuciones para obtener un tiempo promedio en cada uno de los casos.

\end{document}